\section{Úvod}

Dynamika písania (angl.\ \emph{keystroke dynamics}) je behaviorálna biometria,
ktorá z času stlačenia a uvoľnenia kláves odvodzuje charakteristický rytmus
písania konkrétneho človeka. Nepotrebuje špeciálny hardvér, stačí bežná
klávesnica, a dá sa vyhodnocovať počas písania.

Cieľom tejto práce je navrhnúť, implementovať a otestovať jednoduchý systém,
ktorý (i) zbiera časy stlačenia kláves od účastníkov experimentu pri písaní
tej istej vety a (ii) z jedného nového pokusu určí, ktorému z už známych
účastníkov sa rytmus písania najviac podobá. Riešenie pozostáva z webovej
aplikácie na zber dát, serverovej časti s úložiskom a algoritmu na
identifikáciu. Zdrojový kód je dostupný v repozitári~\cite{repo}.

Dokument je členený takto: kapitola~\ref{sec:teoria} zhŕňa teoretické
východiská, kapitola~\ref{sec:navrh} opisuje návrh riešenia a použitú metódu,
kapitola~\ref{sec:impl} implementáciu, kapitola~\ref{sec:test} testovanie
a vyhodnotenie a kapitola~\ref{sec:prirucka} slúži ako príručka na spustenie
a používanie systému.
